\begin{theorem}
  \label{ent:thm:modeq-iff-dvd-sub}
  For integers $n,a,b$, congruence $a\equiv b\pmod n$ is equivalent to $n\mid a-b$.

  \begin{lean}
    theorem modEq_iff_dvd_sub (n a b : ℤ) : a ≡ b [ZMOD n] ↔ n ∣ a - b
  \end{lean}
\end{theorem}

\begin{proof}
  The remainder formulation gives divisibility of $b-a$.
  Negating a divisibility witness reverses the sign of this difference.

  \begin{lean}
    rw [Int.modEq_iff_dvd]
    constructor

    · rintro ⟨k, hk⟩
      exact ⟨-k, by linarith⟩

    · rintro ⟨k, hk⟩
      exact ⟨-k, by linarith⟩
  \end{lean}
\end{proof}

\begin{theorem}
  \label{ent:thm:modeq-iff-emod-eq}
  Integers are congruent modulo $n$ exactly when their remainders on division by $n$ agree.

  \begin{lean}
    theorem modEq_iff_emod_eq (n a b : ℤ) : a ≡ b [ZMOD n] ↔ a % n = b % n
  \end{lean}
\end{theorem}

\begin{proof}
  This is the remainder definition of the congruence relation.

  \begin{lean}
    Iff.rfl
  \end{lean}
\end{proof}

\begin{theorem}
  \label{ent:thm:modeq-refl}
  For all integers $n,a$, one has $a\equiv a\pmod n$.

  \begin{lean}
    theorem modEq_refl (n a : ℤ) : a ≡ a [ZMOD n]
  \end{lean}
\end{theorem}

\begin{proof}
  The remainder of $a$ equals itself.

  \begin{lean}
    rfl
  \end{lean}
\end{proof}

\begin{theorem}
  \label{ent:thm:modeq-symm}
  If $a\equiv b\pmod n$, then $b\equiv a\pmod n$.

  \begin{lean}
    theorem modEq_symm {n a b : ℤ} (h : a ≡ b [ZMOD n]) : b ≡ a [ZMOD n]
  \end{lean}
\end{theorem}

\begin{proof}
  Reverse the equality of remainders.

  \begin{lean}
    Eq.symm h
  \end{lean}
\end{proof}

\begin{theorem}
  \label{ent:thm:modeq-trans}
  If $a\equiv b\pmod n$ and $b\equiv c\pmod n$, then $a\equiv c\pmod n$.

  \begin{lean}
    theorem modEq_trans {n a b c : ℤ} (h : a ≡ b [ZMOD n]) (h' : b ≡ c [ZMOD n]) :
        a ≡ c [ZMOD n]
  \end{lean}
\end{theorem}

\begin{proof}
  Transitivity of equality applies to the three remainders.

  \begin{lean}
    Eq.trans h h'
  \end{lean}
\end{proof}

\begin{theorem}
  \label{ent:thm:modeq-add}
  If $a\equiv b\pmod n$ and $c\equiv d\pmod n$, then $a+c\equiv b+d\pmod n$.

  \begin{lean}
    theorem modEq_add {n a b c d : ℤ} (h : a ≡ b [ZMOD n]) (h' : c ≡ d [ZMOD n]) :
        a + c ≡ b + d [ZMOD n]
  \end{lean}
\end{theorem}

\begin{proof}
  Write $a-b=nu$ and $c-d=nv$.
  The difference of the sums is $n(u+v)$.

  \begin{lean}
    obtain ⟨u, hu⟩ := (modEq_iff_dvd_sub n a b).1 h
    obtain ⟨v, hv⟩ := (modEq_iff_dvd_sub n c d).1 h'

    exact (modEq_iff_dvd_sub n _ _).2 ⟨u + v, by linarith⟩
  \end{lean}
\end{proof}

\begin{theorem}
  \label{ent:thm:modeq-mul}
  Under the same hypotheses, $ac\equiv bd\pmod n$.

  \begin{lean}
    theorem modEq_mul {n a b c d : ℤ} (h : a ≡ b [ZMOD n]) (h' : c ≡ d [ZMOD n]) :
        a * c ≡ b * d [ZMOD n]
  \end{lean}
\end{theorem}

\begin{proof}
  The identity $ac-bd=(a-b)c+b(c-d)$ expresses the difference as $n(uc+bv)$.

  \begin{lean}
    obtain ⟨u, hu⟩ := (modEq_iff_dvd_sub n a b).1 h
    obtain ⟨v, hv⟩ := (modEq_iff_dvd_sub n c d).1 h'

    exact (modEq_iff_dvd_sub n _ _).2 ⟨u * c + b * v, by
      nlinarith [congrArg (fun z => z * c) hu, congrArg (fun z => z * b) hv]⟩
  \end{lean}
\end{proof}

\begin{theorem}
  \label{ent:thm:modeq-pow}
  If $a\equiv b\pmod n$, then $a^k\equiv b^k\pmod n$ for every natural number $k$, including zero.

  \begin{lean}
    theorem modEq_pow {n a b : ℤ} (h : a ≡ b [ZMOD n]) (k : ℕ) :
        a ^ k ≡ b ^ k [ZMOD n]
  \end{lean}
\end{theorem}

\begin{proof}
  At exponent zero both sides are one.

  \begin{lean}
    induction k with
    | zero =>
      simp only [pow_zero]
      exact modEq_refl n 1
  \end{lean}

  Multiply the induction hypothesis by the original congruence to pass to the next exponent.

  \begin{lean}
    | succ k ih => simpa only [pow_succ] using modEq_mul ih h
  \end{lean}
\end{proof}

\begin{theorem}
  \label{ent:thm:modeq-zero-iff}
  For integers $n,a$, one has $a\equiv0\pmod n$ if and only if $n\mid a$.

  \begin{lean}
    theorem modEq_zero_iff (n a : ℤ) : a ≡ 0 [ZMOD n] ↔ n ∣ a
  \end{lean}
\end{theorem}

\begin{proof}
  Set the second integer equal to zero in the divisibility characterization.

  \begin{lean}
    simpa only [sub_zero] using modEq_iff_dvd_sub n a 0
  \end{lean}
\end{proof}

\begin{theorem}
  \label{ent:thm:existsunique-residue}
  For every positive modulus $n$ and integer $a$, there is exactly one natural number $r<n$ with
  $a\equiv r\pmod n$.

  \begin{lean}
    theorem existsUnique_residue (n : ℕ+) (a : ℤ) :
        ∃! r : ℕ, r < n ∧ a ≡ (r : ℤ) [ZMOD (n : ℤ)]
  \end{lean}
\end{theorem}

\begin{proof}
  The division algorithm gives $a=qn+r$ with $0\le r<n$.
  The identity $a-r=qn$ shows that $r$ is a congruent representative.

  \begin{lean}
    obtain ⟨⟨q, r⟩, ⟨ha, hr⟩, hu⟩ := Chapter01.existsUnique_quotient_remainder a n

    refine ⟨r, ⟨hr, ?_⟩, ?_⟩

    · apply (modEq_iff_dvd_sub _ _ _).2
      exact ⟨q, by
          change a = q * (n : ℤ) + r at ha
          nlinarith⟩
  \end{lean}

  Conversely, any congruent representative $s$ in this interval supplies a quotient through
  $a-s=kn$.
  Uniqueness of the quotient--remainder pair forces $s=r$.

  \begin{lean}
    · intro s hs

      obtain ⟨k, hk⟩ := (modEq_iff_dvd_sub _ _ _).1 hs.2

      have hp : (k, s) = (q, r) := hu (k, s) ⟨by
          dsimp
          nlinarith, hs.1⟩

      exact congrArg Prod.snd hp
  \end{lean}
\end{proof}

\begin{definition}
  A family of $n$ integers is a \emph{complete residue system} modulo $n$ if every integer is
  congruent to exactly one member of the family.

  \begin{lean}
    def IsCompleteResidueSystem (n : ℕ) (a : Fin n → ℤ) : Prop :=
      ∀ x : ℤ, ∃! i : Fin n, x ≡ a i [ZMOD (n : ℤ)]
  \end{lean}
\end{definition}

\begin{theorem}
  \label{ent:thm:complete-residue-system-iff}
  For $n>0$, a family $(a_i)_{i\in\{0,\ldots,n-1\}}$ is a complete residue system precisely when
  congruent members have equal indices.

  \begin{lean}
    theorem complete_residue_system_iff (n : ℕ) (hn : 0 < n) (a : Fin n → ℤ) :
        IsCompleteResidueSystem n a ↔ ∀ i j : Fin n, a i ≡ a j [ZMOD (n : ℤ)] → i = j
  \end{lean}
\end{theorem}

\begin{proof}
  Uniqueness of a representative immediately gives pairwise incongruence.

  \begin{lean}
    let : NeZero n := ⟨ne_of_gt hn⟩

    constructor

    · intro h i j hij

      obtain ⟨k, hk, hu⟩ := h (a i)

      exact (hu i (modEq_refl _ _)).trans (hu j hij).symm
  \end{lean}

  Conversely, pairwise incongruence makes the map from the $n$ indices to the $n$ residue classes
  injective.

  \begin{lean}
    · intro h

      have hinj : Function.Injective (fun i : Fin n => (a i : ZMod n)) := by
        intro i j heq
        exact h i j ((ZMod.intCast_eq_intCast_iff _ _ n).1 heq)
  \end{lean}

  Finite counting makes this map surjective as well.
  Every integer therefore has a representative, and injectivity makes it unique.

  \begin{lean}
    have hsurj := (Finite.injective_iff_surjective_of_equiv (ZMod.finEquiv n).toEquiv).1 hinj

    intro x

    obtain ⟨i, hi⟩ := hsurj (x : ZMod n)

    refine ⟨i, (ZMod.intCast_eq_intCast_iff _ _ n).1 hi.symm, ?_⟩

    intro j hj
    apply hinj
    exact ((ZMod.intCast_eq_intCast_iff _ _ n).2 hj).symm.trans hi.symm
  \end{lean}
\end{proof}

\begin{theorem}
  \label{ent:thm:modeq-cancel}
  If $\gcd(n,c)=1$ and $ca\equiv cb\pmod n$, then $a\equiv b\pmod n$.

  \begin{lean}
    theorem modEq_cancel {n c a b : ℤ} (hc : Chapter01.gcd n c = 1)
        (h : c * a ≡ c * b [ZMOD n]) : a ≡ b [ZMOD n]
  \end{lean}
\end{theorem}

\begin{proof}
  The hypothesis says $n\mid c(a-b)$.
  The previously proved coprime form of Euclid\textquotesingle s lemma gives $n\mid a-b$.

  \begin{lean}
    apply (modEq_iff_dvd_sub n a b).2
    apply Chapter01.dvd_of_dvd_mul_of_coprime n c (a - b) hc
    simpa only [mul_sub] using (modEq_iff_dvd_sub n _ _).1 h
  \end{lean}
\end{proof}

\begin{theorem}
  \label{ent:thm:modeq-cancel-gcd}
  Let $n\ne0$ and $g=\gcd(n,c)$.
  If $ca\equiv cb\pmod n$, then $a\equiv b\pmod{n/g}$.

  \begin{lean}
    theorem modEq_cancel_gcd {n c a b : ℤ} (hn : n ≠ 0)
        (h : c * a ≡ c * b [ZMOD n]) :
        a ≡ b [ZMOD (n / (Chapter01.gcd n c : ℤ))]
  \end{lean}
\end{theorem}

\begin{proof}
  Write $n=gu$ and $c=gv$.
  Dividing the pair by its positive gcd gives $\gcd(u,v)=1$.

  \begin{lean}
    obtain ⟨hd, ⟨u, hu⟩, ⟨v, hv⟩, _, _⟩ := Chapter01.gcd_spec n c (Or.inl hn)

    have hcop := Chapter01.gcd_div_gcd_eq_one n c (Or.inl hn)
    generalize hg : Chapter01.gcd n c = g at *

    have hdne : (g : ℤ) ≠ 0 := by
      exact_mod_cast ne_of_gt hd

    have hnu : n / (g : ℤ) = u := by
      rw [hu]
      exact Int.mul_ediv_cancel_left u hdne

    have hcv : c / (g : ℤ) = v := by
      rw [hv]
      exact Int.mul_ediv_cancel_left v hdne
    rw [hnu, hcv] at hcop
  \end{lean}

  A witness for $n\mid c(a-b)$ yields $gu k=gv(a-b)$.
  Cancel the nonzero factor $g$, then apply Euclid\textquotesingle s lemma to $u\mid v(a-b)$.

  \begin{lean}
    obtain ⟨k, hk⟩ := (modEq_iff_dvd_sub n _ _).1 h

    have hdiv : u ∣ v * (a - b) := by
      refine ⟨k, ?_⟩
      apply mul_left_cancel₀ hdne
      calc
        (g : ℤ) * (v * (a - b)) = c * a - c * b := by
          rw [hv]
          ring
        _ = (g : ℤ) * (u * k) := by
          rw [hk, hu]
          ring
    rw [hnu]
    exact (modEq_iff_dvd_sub _ _ _).2
      (Chapter01.dvd_of_dvd_mul_of_coprime u v (a - b) hcop hdiv)
  \end{lean}
\end{proof}

\begin{theorem}
  \label{ent:thm:modeq-mul-zero-prime}
  For a prime $p$, the congruence $ab\equiv0\pmod p$ implies $a\equiv0\pmod p$ or $b\equiv0\pmod
    p$.

  \begin{lean}
    theorem modEq_mul_zero_prime (p : ℕ) (hp : Chapter02.isPrime p) (a b : ℤ)
        (h : a * b ≡ 0 [ZMOD (p : ℤ)]) :
        a ≡ 0 [ZMOD (p : ℤ)] ∨ b ≡ 0 [ZMOD (p : ℤ)]
  \end{lean}
\end{theorem}

\begin{proof}
  Translate the congruences into divisibility and apply the prime-divides-a-product theorem from
  Chapter 2.

  \begin{lean}
    rw [modEq_zero_iff] at h ⊢
    rw [modEq_zero_iff]
    exact Chapter02.prime_dvd_mul p hp a b h
  \end{lean}
\end{proof}
